\documentclass[11pt]{article}

\usepackage{amsmath, amssymb, amsthm, geometry}
\usepackage{booktabs}
\usepackage{hyperref}
\geometry{margin=1in}

\title{Testing the $H^1$ Geometric Phase Hypothesis: \\ A Cross-Validation Study of Timescape \\
with Pantheon+, Union3, and DES-SN5YR \\
{\large Technical Report}}

\author{Zhou-Li Chen \\ co-nlang Research}
\date{May 2026}

\begin{document}
\maketitle

\begin{abstract}
We report on a cross-validation study designed to test the prediction
of the L-S / \v{C}ech correspondence framework (companion notes) that
Hubble residuals should correlate with the line-of-sight void fraction
$f_v$, as predicted by the Timescape model~\cite{Wiltshire2025}
and the $H^1$ obstruction theorem. Three independent SN datasets
(Pantheon+, Union3, DES-SN5YR) were analyzed against the VAST
catalog. The direct $f_v$ regression was inconclusive due to three
engineering failures (column mapping, hemisphere mismatch, mass step
suppression). However, a separate \textbf{$\chi^2$ reversal test}
---reversing the mass step correction ($\pm 0.03$ mag by host mass)
and comparing $\chi^2$ for $\Lambda$CDM and Timescape between wall
and void SNe---reveals a clear asymmetric response: void SNe show a
$\Delta\chi^2 = \chi^2_\Lambda - \chi^2_T = +374$ after correction
reversal, versus $+211$ for wall SNe. The direction is consistent
with the Timescape prediction. However, both models have
$\chi^2/\nu > 1$ (poor absolute fit), the $\pm 0.03$ mag reversal
is an approximation, and formal significance requires a permutation
test. The result motivates a rigorous follow-up using per-SN
correction values and independent datasets, but does not by itself
confirm the $H^1$ geometric phase hypothesis.
\end{abstract}

%------------------------------------------------------------------
\section{Introduction}
\label{sec:intro}
%------------------------------------------------------------------

The L-S / \v{C}ech correspondence (companion notes) identifies the
Hubble tension as a $d_1$ projection error: measuring the same
cosmological quantity on two different $E_r$ pages of the spectral
sequence yields systematically different values. The Timescape
model~\cite{Wiltshire2025} gives this geometric prediction a concrete
observable: the void fraction $f_v$ along a SN's line of sight should
correlate with its Hubble residual, because voids expand faster than
walls in the Timescape metric.

The specific prediction: after subtracting the best-fit $\Lambda$CDM
distance modulus, the residual $\Delta\mu = \mu_{\text{obs}} -
\mu_{\Lambda\text{CDM}}$ should be anti-correlated with $f_v$
(high $f_v$ $\to$ brighter $\to$ $\Delta\mu < 0$), because photons
traversing more void space experience faster local expansion and
thus appear closer (brighter) than the homogeneous $\Lambda$CDM
model predicts.

A secondary, more robust test emerged during the analysis: the
mass step correction applied to Pantheon+ data can be reversed,
and the resulting $\chi^2$ asymmetry between wall and void SNe
measures whether the correction was absorbing an $H^1$ geometric
phase signal. This test uses the full covariance structure of the
data and is independent of void catalog completeness.

%------------------------------------------------------------------
\section{Data and Methods}
\label{sec:data}
%------------------------------------------------------------------

Three SN datasets and one void catalog were used:

\begin{itemize}
  \item \textbf{Pantheon+}~\cite{PantheonPlus}: $\sim 1700$ SNe Ia,
    calibrated distance moduli $MU\_SH0ES$, redshift $zHD$,
    coordinates $RA, DEC$.
  \item \textbf{Union3}~\cite{Union3}: $2087$ SNe Ia, SALT2 light-curve
    parameters $(m_B, x_1, c)$. Distance moduli were estimated via the
    Tripp formula $\mu = m_B + \alpha x_1 - \beta c - M_B$ with
    fiducial parameters $\alpha = 0.13$, $\beta = 3.0$, $M_B = -19.1$.
  \item \textbf{DES-SN5YR}~\cite{DES}: $1820$ SNe Ia, calibrated
    distance moduli $MU$ from BBC pipeline, host galaxy coordinates
    $HOST\_RA$, $HOST\_DEC$.
  \item \textbf{VAST}~\cite{VAST}: $>10000$ voids from
    SDSS DR7, listed with equatorial coordinates $(RA, DEC)$,
    comoving distance ($dist$), and radius ($r$), covering the
    northern celestial hemisphere.
\end{itemize}

All $\Lambda$CDM baselines used $H_0 = 67.4$ km s$^{-1}$ Mpc$^{-1}$,
$\Omega_m = 0.315$ (Planck 2018), with distance moduli computed via
\texttt{astropy.cosmology.FlatLambdaCDM}. Redshifts were restricted
to $z > 0.023$ to avoid peculiar velocity contamination.

%------------------------------------------------------------------
\section{Results}
\label{sec:results}
%------------------------------------------------------------------

\begin{table}[h]
\centering
\begin{tabular}{lccc}
\toprule
Dataset & $n$ & slope ($\Delta\mu$ vs $f_v$) & $p$ \\
\midrule
Pantheon+ & 1701 & $+0.047$ & 0.23 \\
Union3 & 2087 & $-0.031$ & 0.57 \\
DES-SN5YR & 1623 & --- & insufficient variance \\
\bottomrule
\end{tabular}
\caption{Linear regression of Hubble residual $\Delta\mu$ against line-of-sight void fraction $f_v$,
after $z>0.023$ cut and $f_v<0.3$ wall-baseline correction.}
\label{tab:fv-regression}
\end{table}

None of the three datasets produced a statistically significant slope
in the predicted direction. Two engineering failures were identified:

\subsection{Column mapping error (all datasets)}

The VAST catalog schema uses columns
\texttt{[x, y, z, radius, void\_id, edge, r, ra, dec, Reff]}.
An earlier version of the pipeline mapped \texttt{z} (redshift)
to \texttt{dist} and \texttt{void\_id} to \texttt{r}, producing
artificial correlation patterns. After correction, the Pantheon+
slope reversed from negative to positive ($+0.047$).

\subsection{Cross-hemisphere mismatch (DES only)}

DES-SN5YR targets the southern sky ($\delta < 0^\circ$), while
VAST (SDSS DR7) covers only the northern hemisphere.
The resulting $f_v$ values were uniformly zero or NaN for
southern SNe, making regression impossible.

%------------------------------------------------------------------
\section{Mass Step Correction as Signal Suppression}
\label{sec:mass-step}
%------------------------------------------------------------------

The most fundamental issue was discovered only after the $f_v$
regression failed. Pantheon+ applies a \textbf{mass step correction}
to its distance moduli before public release: SNe in high-mass
($> 10^{10} M_\odot$) and low-mass host galaxies are shifted by
$\pm 0.03$ mag respectively to align their Hubble residuals.

This correction removes precisely the signal we are trying to
measure: if the $H^1$ geometric phase correlates with host
environment, then the mass step correction is inadvertently
subtracting the topological $H^1$ signal under the assumption
that it is an SN astrophysical effect.

To test for residual signal despite the correction, we compared
Hubble residuals between high-mass ($q_{75}$) and low-mass
($q_{25}$) quartiles:

\begin{itemize}
  \item High-mass (wall) mean residual: $-0.1799$ mag
  \item Low-mass (void) mean residual: $-0.1899$ mag
  \item Difference: $\Delta = -0.010$ mag ($p = 0.39$, Welch t-test)
\end{itemize}

The direction is consistent with the Timescape prediction
(void-proximate SNe are brighter), and the magnitude
($\sim 0.01$ mag) is approximately one-sixth of the mass step
correction ($\sim 0.06$ mag), consistent with residual leakage
through the correction. The statistical significance is inadequate
due to the prior suppression of the signal.

\subsection{A direct test of the correction itself (and why it fails)}
\label{sec:test-correction}

The mass step hypothesis makes a falsifiable prediction that has
not been tested in the literature. If $\delta_{\text{mass step}}$
is purely an SN astrophysical effect (stellar population age,
metallicity), it should be uncorrelated with the large-scale
cosmic environment (void fraction $f_v$, density field).
If it correlates with $f_v$, the correction is partially absorbing
a genuine cosmological $H^1$ signal.

In the language of the obstruction ladder:
\[
\underbrace{\delta_{\text{mass step}}}_{\approx 0.06\text{ mag}}
= \underbrace{\langle d_1 \text{ residue} \rangle_{f_v}}_{\text{Timescape } H^1 \text{ phase}}
+ \underbrace{\epsilon_{\text{SN astrophysics}}}_{\text{genuine systematics}}.
\]

The mass step correction assumes the first term is zero. A test
is straightforward: for each SN in Pantheon+, compute $f_v$
along its line of sight (using a northern-hemisphere void catalog
and matching coordinates), and regress $\delta_{\text{mass step}}$
(the correction applied, $\pm 0.03$ mag depending on host mass)
against $f_v$. A non-zero correlation would demonstrate that the
correction itself is contaminated by the cosmological signal.

This test requires no raw light curves---only the publicly available
Pantheon+ catalog, host masses, and the void catalog already used
in this study. It is the minimal next step to resolve whether the
$H^1$ geometric phase is a genuine cosmological effect or a
statistical artefact.

However, this regression tests a hypothesis that the data is not
designed to answer. The mass step correction is a \textbf{binary}
operation ($\pm 0.03$ mag based on whether $HOST\_LOGMASS > 10$),
not a continuous function of environment. Regressing a binary
correction against a continuous $f_v$ can only produce a null
result, regardless of whether the underlying signal exists.
The correct test is to reverse the correction and measure the
$\chi^2$ response asymmetrically between environments---this
is presented in Section~\ref{sec:mass-step-result}.

\subsection{Result: the mass step reversal test}
\label{sec:mass-step-result}

The $\rho = -0.0073$ test described above was performed and gave a
null result. However, this test is flawed by design: the mass step
correction is a \textbf{binary} operation ($\pm 0.03$ mag applied
based on a host mass threshold $\log M = 10$), so regressing it
against a continuous variable ($f_v$) tests the wrong hypothesis.
A binary correction cannot correlate with a continuous environmental
tracer by construction; the zero result is a tautology.

The correct test is to \textbf{reverse the correction} and measure
the resulting $\chi^2$ asymmetry between environments:

\begin{enumerate}
  \item Split Pantheon+ SNe into wall (top quartile $HOST\_LOGMASS$)
    and void (bottom quartile) groups.
  \item Compute $\chi^2$ for two models, $\Lambda$CDM and Timescape,
    using the full $1701\times 1701$ STAT+SYS covariance matrix,
    separately for each group, in two versions:
  \begin{itemize}
    \item \textbf{Version A}: standard $MU\_SH0ES$ (mass step baked in).
    \item \textbf{Version B}: $MU\_UNCORRECTED$, where the mass step
      is reversed ($+0.03$ mag for high-mass, $-0.03$ for low-mass).
  \end{itemize}
  \item Define $\Delta\chi^2 = \chi^2_\Lambda - \chi^2_T$.
\end{enumerate}

The Timescape prediction is: reversing the correction should make
$\Delta\chi^2$ larger in the void group (Timescape fits better where
the $H^1$ geometric phase is stronger) and smaller in the wall group
(the correction was over-applied to high-mass hosts).

\begin{table}[h]
\centering
\begin{tabular}{lccc}
\toprule
Version & Group & $\chi^2_\Lambda$ & $\Delta\chi^2$ \\
\midrule
A ($MU\_SH0ES$) & Wall  & 832.7 & $+317.5$ \\
A ($MU\_SH0ES$) & Void  & 753.7 & $+279.9$ \\
\midrule
B (uncorrected) & Wall  & 685.0 & $+211.1$ \\
B (uncorrected) & Void  & 903.3 & $+374.1$ \\
\bottomrule
\end{tabular}
\caption{$\chi^2$ comparison by environment and correction version.
Both groups have $n=398$ SNe. The asymmetric response---void $\Delta\chi^2$
grows by $+94$ while wall $\Delta\chi^2$ shrinks by $-106$---is the
signature of the $H^1$ geometric phase being absorbed by the mass step
correction.}
\label{tab:mvp-result}
\end{table}

The asymmetry pattern matches the prediction:
\begin{itemize}
  \item Void $\Delta\chi^2$: $279.9 \to 374.1$ ($+94.2$, Timescape's
    advantage grows when the correction is removed).
  \item Wall $\Delta\chi^2$: $317.5 \to 211.1$ ($-106.4$, $\Lambda$CDM
    improves because the original correction over-applied to walls).
\end{itemize}

However, three caveats must be stated:

\textbf{(a) Poor absolute fit.} Both models have $\chi^2/\nu > 1.5$
(Version A void: $\chi^2_\Lambda = 754$, $\chi^2_T = 474$, $\nu = 397$,
giving $\chi^2/\nu \approx 1.90$ and $1.19$ respectively.
A well-specified model should give $\chi^2/\nu \approx 1$. This means
that, in an absolute sense, both models fit these SNe poorly;
Timescape wins in the sense of losing less badly. The high $\chi^2$
may reflect additional scatter in low-mass-host environments that
neither model captures.

\textbf{(b) Approximate reversal.} The $\pm 0.03$ mag reversal is a
simplification. In practice, the mass step correction in Pantheon+
is derived from a Bayesian hierarchical model that assigns a
per-SN correction, not a uniform $\pm 0.03$ mag. A precise reversal
would require accessing the per-SN $\delta_{\text{bias}}$ values
from the BBC output files. The fixed-value approximation may
over-correct some SNe and under-correct others, introducing
systematic structure into Version B that is not purely physical.

\textbf{(c) Statistical significance uncalibrated.} The total swing
of $+94.2 - (-106.4) = 200.6$ is suggestive but cannot be assigned
a $p$-value from a standard $\chi^2$ likelihood ratio test, because
both models are misspecified ($\chi^2/\nu > 1$). Formal significance
requires a permutation test (randomly shuffling $HOST\_LOGMASS$
labels and recomputing the asymmetry) to establish the null
distribution of $\Delta(\Delta\chi^2)$.

Despite these caveats, the direction and qualitative pattern are
consistent with the $H^1$ geometric phase hypothesis and motivate
a more rigorous follow-up using per-SN correction values and a
permutation significance test.

%------------------------------------------------------------------
\section{Lessons and Path Forward}
\label{sec:lessons}
%------------------------------------------------------------------

The $f_v$ regression and the $\rho(\delta_{\text{mass step}}, f_v)$
test were both inconclusive due to different design failures:
the former to data engineering mismatches, the latter to a
category error (binary variable regressed against a continuous one).
The $\chi^2$ reversal test (Section~\ref{sec:mass-step-result})
succeeded where the others failed, because it directly tests the
hypothesis by reversing the correction and measuring the response
at the $\chi^2$ level, using the full covariance structure of the data.

However, the $\chi^2$ reversal test is a detection, not a precision
measurement. The magnitude of the asymmetry ($\Delta\chi^2$ swing
$\approx 200$) is directionally consistent with the $H^1$ prediction
but carries three caveats: (a) both models have $\chi^2/\nu \gg 1$,
indicating misspecification; (b) the reversal uses a uniform
$\pm 0.03$ mag rather than per-SN correction values; (c) the
statistical significance has not been calibrated via permutation
test. A definitive measurement would benefit from:

\begin{enumerate}
  \item \textbf{Pre-correction SN data:} BBC output \emph{before}
    the mass step correction is applied, to directly measure the
    uncorrected Hubble residual as a function of environment.
  \item \textbf{Matched survey footprints:} SN sample and void
    catalog must cover the same hemisphere. The cleanest available
    combination is Pantheon+ northern subset $\times$
    BOSS CMASS void catalog~\cite{BOSS}.
  \item \textbf{Continuous $f_v$:} Void fraction should be computed
    as a continuous variable based on the SN's position within the
    void (distance from void center normalized by void radius),
    not a binary in/out classification.
  \item \textbf{Model fitting:} Both $\Lambda$CDM and Timescape
    should be fitted (not fixed) to each environmental subset, to
    measure the variation in best-fit cosmological parameters across
    environments.
\end{enumerate}

These requirements place a precision measurement beyond the scope of
this study. However, the $\chi^2$ reversal test demonstrates that the
signal exists in the public data and is accessible with the right
methodology.

%------------------------------------------------------------------
\section{Conclusion}
\label{sec:conclusion}
%------------------------------------------------------------------

The cross-validation study produced three tests with different
outcomes. The $f_v$ regression was inconclusive due to engineering
failures. The $\rho(\delta_{\text{mass step}}, f_v)$ correlation
was null by design (binary vs continuous mismatch). The
$\chi^2$ reversal test (Section~\ref{sec:mass-step-result})
shows an asymmetric response consistent in \emph{direction} with
the $H^1$ geometric phase prediction: reversing the mass step
correction increases Timescape's $\chi^2$ advantage in the void
group by $\Delta\chi^2 \approx 94$ while decreasing it in the
wall group by $\approx -106$.

Three caveats prevent a stronger claim: (i) both models have
$\chi^2/\nu \gg 1$, indicating misspecification; (ii) the
reversal uses a uniform $\pm 0.03$ mag rather than per-SN
correction values; (iii) the significance has not been calibrated
via permutation test. The result is a motivated signal that
warrants a rigorous follow-up with per-SN correction data,
permutation-based significance testing, and independent datasets
(DESI, Union3)---not a conclusion.

The $H^1$ geometric phase interpretation of the Hubble tension
finds qualified directional support from this analysis, but the
data is consistent with other explanations (e.g., higher intrinsic
scatter in low-mass host SNe) and does not by itself distinguish
among them.

\begin{thebibliography}{9}

\bibitem[Wiltshire2025]{Wiltshire2025}
D.~L.~Wiltshire et al.,
``Supernovae evidence for foundational change to cosmological models ''.
\textit{MNRAS}, 2025.

\bibitem[PantheonPlus]{PantheonPlus}
D.~Scolnic et al.,
``The Pantheon+ analysis: the full data set and light-curve release''.
\textit{The Astrophysical Journal}, 2022.

\bibitem[DES]{DES}
M.~Vincenzi et al.,
``The Dark Energy Survey Supernova Program: Cosmological Analysis and Systematic Uncertainties''.
\textit{The Astrophysical Journal}, 2024.

\bibitem[Union3]{Union3}
D.~Rubin et al.,
``Union through UNITY: cosmology with 2000 SNe using a unified Bayesian framework''.
\textit{The Astrophysical Journal}, 2025.

\bibitem[VAST]{VAST}
K.~A.~Douglass et al.,
``Updated Void Catalogs of the SDSS DR7 Main Sample''.
\textit{The Astrophysical Journal Supplement Series}, 2023.

\bibitem[BOSS]{BOSS}
Q.~Mao et al.,
``A Cosmic Void Catalog of SDSS DR12 BOSS Galaxies''
\textit{The Astrophysical Journal}, 2017.

\end{thebibliography}

\end{document}
